\section{Compilation over the original alphabet}
\label{sec:compilation}

The semantic construction is complete: three splits and Boolean
recoveries determine the monoid value on episode sequences. To turn this
into a height-four expression, we must account for the component languages
and the final remainder. This section proves both claims directly over
$\Sigma$. In particular, a suffix language must have its own height-one
expression even when its guessed origin is wrong and no matching prefix
exists.

The common mechanism is a finite collection of literal word templates
$a b^h c$. Every component has at most one unbounded continuation; it
follows one fixed period until its first mismatch. On a fixed template,
all other tests become ultimately periodic conditions on the exponent
$h$. This provides the expression itself, rather than inferring a height
bound from mere regularity.

\subsection{A template lemma}

\begin{lemma}[Periodic templates]
\label{lem:templates}
Consider a finite collection of templates
\begin{equation}
 \label{eq:periodic-template}
 a b^h c,\qquad h\ge0,
\end{equation}
where $a,b,c\in\Sigma^*$ are fixed for each template and $b\ne\eps$.
Suppose a test on each template uses only the following data:
\begin{enumerate}[label=\textup{(\roman*)}]
\item letters in fixed finite collections of windows at offsets from
the argument's start or end, with mandatory availability checks;
\item order and availability of endpoints chosen from fixed finite sets
of offsets from the start or end;
\item $M$-products between these endpoints, whenever available and
ordered, and lengths modulo fixed positive integers;
\item fixed finite combinations, choices, and calculations from these
data.
\end{enumerate}
The accepted words have a generalized expression of height at most one.
Adding or removing a finite language preserves this bound.

Moreover, any of the exact-end tests of Section~\ref{sec:episodes},
with its reference or decoded origin chosen from finitely many offsets
relative to the test start, is covered by a finite language and finitely
many templates of the form~\eqref{eq:periodic-template}.
\end{lemma}

\begin{proof}
\emph{Ultimately periodic data and explicit expressions.}
Fix a template.  For sufficiently large $h$, all the indicated
availability and order relations stabilize.  Letters at a bounded
distance from the start or end also stabilize, since increasing $h$
inserts a full copy of $b$.  A product between boundedly separated
endpoints then stabilizes.  A product spanning the growing middle has
the form
\[
 u\,T(b)^{h-j}\,v
\]
for fixed $u,v\in M$ and a fixed integer $j$, once $h$ is sufficiently
large.  This follows by separating its fixed initial and terminal pieces
from its complete copies of $b$.  Powers of an element of a finite monoid
are ultimately periodic: if two powers coincide, multiplying by further
powers gives the eventual period.  Lengths modulo any fixed integer are
periodic in $h$.  Thus the entire finite data array is ultimately
periodic.  A fixed finite computation from this array, including a
universal or uniqueness check on a fixed finite list, remains ultimately
periodic.

Let the resulting accepted exponents have period $f>0$ after a threshold.
They are a finite set together with finitely many progressions
$h_0+f\mathbb N$.  Each progression contributes the expression
\[
 a\,b^{h_0}\,(b^f)^*\,c,
\]
where the displayed fixed powers are finite words.  These expressions
have height one; the exceptional words have height zero.  A finite union
proves the first assertion, including a finite number of templates and
finite exceptions.

\emph{The template cover for an independent exact-end test.}
For the final assertion, fix one possible end-rule origin relative to the
argument start.  The present-marker branch has a bounded anchor and
therefore bounded total length.  For the absent-marker branch, the seed
is at a fixed offset and has fixed length.  Make finite choices of all
letters through the seed, of its short period and phase, and of the
fixed-length tail beginning at the first mismatch.  These choices fix
a prefix $a$ through the seed and a nonempty period word $b$ beginning
at the next position.  If the mismatch occurs after $h$ full copies of
$b$ and a partial copy of length $r<|b|$, place that partial copy, the
mismatching letter, and its remaining $b_{\rm tail}-1$ letters in $c$.
There are finitely many such choices, and the mismatching letter is
required to differ from the next periodic letter.  They give exactly
templates~\eqref{eq:periodic-template} for this continuation and end.
The bounded marker and availability conditions can still be tested on
the templates.  Taking the finite union over possible end-rule origins
proves the assertion.
\end{proof}

This proof applies to the exact-end rule of a suffix considered by itself.
It does not presuppose that its guessed origin is correct, or that it is
actually the suffix of an episode.  If a hypothesized origin lacks a
mandatory seed or window, it is rejected.  Otherwise its own exact-end
rule supplies the template.

\begin{proposition}[Component bounds]
\label{prop:components}
The episode languages $E,D_0,D_1$ and every suffix component $Q_y$ in
Section~\ref{sec:splits} have generalized star height at most one.
Every prefix component $P_x$ has height zero.  These bounds hold over
the original alphabet $\Sigma$.
\end{proposition}

\begin{proof}
We first identify the single end rule that supplies each template cover,
then verify that every remaining operation is a finite calculation from
the allowed data. No part of this argument assumes that a standalone
suffix has guessed the true episode origin.

All prefix cuts come from fixed finite sets of offsets, so the lengths
of $P_x$ words are bounded.  Hence $P_x$ is finite.

Every other language in the assertion has one exact-end rule on its
entire argument.  For $E,D_0,D_1$ this is the true episode rule.  A main
or stage-0 suffix uses the reference rule in its end guard.  A stage-1
suffix uses the rule at its uniquely decoded candidate. For a stage-1
suffix considered alone, first make a finite case distinction for the
decoded offset $p'$. The uniqueness predicate inspects all candidates'
bounded end-search windows and their finitely many $\gamma$ translates;
it contains no mismatch search. The case then has just the exact-end
rule at $k+L-p'$, whether or not that is a genuine episode origin.
There are only
finitely many reference or decoded offsets; we may take a finite union
over them.  Lemma~\ref{lem:templates} therefore supplies the finite
template cover even for standalone suffixes with incorrect hypotheses.

We check that the remaining tests use only the finite data in that lemma.
Marker searches and their stability checks inspect bounded windows at
fixed nominal offsets.  So do usability, the phase and lattice choices,
and all additional candidate translates.  Every selected main boundary
lies within a fixed bounded neighborhood of its center.  Buffer
boundaries have fixed offsets.  Consequently all partial products in
the component computations have endpoints selected from finite sets of
start offsets, except for the argument end itself.

The anchor supplied by the exact-end rule is the argument end minus
$b_{\rm tail}$.  Thus the clock and $A_0$ use only fixed-modulus
differences between this end offset and their candidate origins.
The predicates $A_1$ and the stage-1 decoding predicate are wholly
bounded.  Equality, universal checks, and uniqueness over their finite
candidate lists are finite calculations, with the mandatory availability
convention retained for every candidate.

The affine-line checks range over finite vector spaces and the finite
monoid $M$.  Their data consist of the specified states and actual
partial products just described.  The tables $\Phi_i$ and $H_i$ are
computed by their explicit formulas using hypothetical elements of $M$
and actual suffix products and schedule data.  Their finite recursion
does not invoke a graph test on an episode sequence.  The total-table
property permits every hypothetical monoid value, so no extra prefix
existence test is being imposed.

Finally, no translated candidate makes a second unbounded continuation
search.  The guards compare bounded marker data and use just their one
reference rule.  Stage~1 makes bounded decoding checks before applying
the one decoded rule.  Accordingly all remaining acceptance conditions
are fixed finite calculations from the template data.  Applying
Lemma~\ref{lem:templates} proves the bounds.  Every resulting expression
uses finite words in $\Sigma^*$ and Boolean operations, concatenation,
and stars over $\Sigma$; no auxiliary tag has become an input letter.
\end{proof}

\subsection{The final remainder}

The episode rule does not require a word to end in a complete episode.
Define
\[
 F=\neg(E\Sigma^*),\qquad \Sigma^*=\neg0.
\]
Thus $F$ consists of words with no episode prefix and has height at most
one.  To finish, we need the same bound for its monoid-value tests.

\begin{lemma}[Remainder]
\label{lem:remainder}
For every $d\in M$, the language $F\cap T^{-1}(d)$ has height at most
one.  Every word admits a unique factorization into a sequence of
episodes followed by a word in $F$.
\end{lemma}

\begin{proof}
\emph{Product tests on the incomplete tail.}
Choose a fixed length threshold exceeding all the bounded decision
windows at a first start and all possible first-episode lengths in the
present-marker case.  A longer word in $F$ cannot have a present end
marker: that would already give an episode prefix.  Its marker-absent
seed has a fixed short-period continuation.  Either no mismatch has yet
occurred, or its first mismatch is at a position $q$ satisfying
\[
 |u|-q<b_{\rm tail}
\]
when the word $u$ starts at position zero.  Otherwise the prefix ending
at $q+b_{\rm tail}$ is a complete episode.  All the required bounded
windows for that prefix are internal by Lemma~\ref{lem:placement},
even when the mismatch itself is early.

In the no-mismatch case, finite choices of the bounded initial word,
period, phase, and final partial period give templates $a b^h c$.
In the mismatch case, append one of the finitely many possible
post-mismatch portions of length less than $b_{\rm tail}$; these again
give finitely many such templates.  Words below the threshold form a
finite language.  On every template the test $T(u)=d$ is an ultimately
periodic monoid-product test, so Lemma~\ref{lem:templates} gives a
height-one language with that value on the template cover.  Intersecting
with $F$ gives exactly $F\cap T^{-1}(d)$ and keeps the same bound.

\emph{Unique parsing.}
Repeatedly remove an episode prefix while one exists.  Every episode is
nonempty, so this procedure terminates, and the remaining word is in
$F$.  Lemma~\ref{lem:prefix-code} makes each removed prefix unique.
Any alternative $E^*F$ factorization must therefore remove the same
successive episodes; it also cannot stop earlier while an episode
prefix remains.  This proves uniqueness.
\end{proof}

\subsection{The height bound}

\begin{proof}[Proof of Theorem~\ref{thm:main}]
The component bounds needed for the three applications of
Lemma~\ref{lem:split} are supplied by
Proposition~\ref{prop:components}.  If the skip graph tests have height
at most $b$, that lemma gives the new bound
\[
 \max\{2,b+1\}.
\]
Indeed the episode-domain test $D^*$ has height at most two, and the
single star in the split formula lies over component and skip tests of
height at most $\max\{1,b\}$.

The stage-1 skipped class is $D_2=\varnothing$, whose graph tests consist
only of the empty-sequence tests and have height zero.  In the order
proved in Proposition~\ref{prop:aligned-splits} and
Proposition~\ref{prop:main-split}, the bounds are therefore
\[
 \underbrace{0}_{D_2^*}
 \ \longrightarrow\ 
 \underbrace{2}_{\text{stage 1}}
 \ \longrightarrow\ 
 \underbrace{3}_{\text{stage 0}}
 \ \longrightarrow\ 
 \underbrace{4}_{\text{main split}}.
\]
The affine linear-part recoveries and distinguished-basis queries use
only finite Boolean operations and do not raise these bounds.  The
linear part of the main episode update is the action of $T_e$ on $V$.
Thus its recovered sequence tests determine $T$ on $E^*$, since the
image of the identity basis vector identifies the monoid value.

Let $A_d=E^*\cap T^{-1}(d)$ be the resulting height-at-most-four tests,
and let $F_d=F\cap T^{-1}(d)$ be the height-at-most-one tests from
Lemma~\ref{lem:remainder}.  For a language recognized by $T$ with
accepting subset $S\subseteq M$, its expression is
\[
 \bigcup_{\substack{d,f\in M\\df\in S}} A_d F_f.
\]
The factorization lemma and the morphism identity prove that this is
exactly the desired language.  Concatenation and finite union do not
increase the maximum height, which is four.  All expressions are over
$\Sigma$, proving the theorem.
\end{proof}
